\documentclass[11pt]{article}

% ---------------------------------------------------------------- packages
\usepackage[T1]{fontenc}
\usepackage{textcomp}
\usepackage{newunicodechar}
\usepackage[utf8]{inputenc}
\usepackage{lmodern}
\usepackage[margin=1in]{geometry}
\usepackage{graphicx}
\usepackage{booktabs}
\usepackage{longtable}
\usepackage{array}
\usepackage{ragged2e}
\usepackage{enumitem}
\usepackage{xcolor}
\usepackage{titlesec}
\usepackage{fancyhdr}
\usepackage[hidelinks]{hyperref}
\usepackage{parskip}

% Hawaiian ʻokina and other diacritics appear in quoted source text; declare
% fallbacks so a quoted sentence can never break the build.
\providecommand{\okina}{'}
\newunicodechar{ʻ}{\okina}
\newunicodechar{ʼ}{'}
\newunicodechar{’}{'}
\newunicodechar{‘}{`}
\newunicodechar{“}{``}
\newunicodechar{”}{''}
\newunicodechar{–}{--}
\newunicodechar{—}{---}
\newunicodechar{…}{\ldots}
\newunicodechar{ō}{\=o}
\newunicodechar{ā}{\=a}
\newunicodechar{ē}{\=e}
\newunicodechar{ī}{\=\i}
\newunicodechar{ū}{\=u}
\newunicodechar{ō}{\=o}
\newunicodechar{≈}{\ensuremath{\approx}}
\newunicodechar{≤}{\ensuremath{\leq}}
\newunicodechar{≥}{\ensuremath{\geq}}
\newunicodechar{→}{\ensuremath{\rightarrow}}
\newunicodechar{·}{\textperiodcentered}
\newunicodechar{∼}{\ensuremath{\sim}}
\newunicodechar{°}{\textdegree}
\newunicodechar{×}{\ensuremath{\times}}
\newunicodechar{µ}{\textmu}
\newunicodechar{μ}{\ensuremath{\mu}}
\newunicodechar{²}{\textsuperscript{2}}
\newunicodechar{³}{\textsuperscript{3}}
\newunicodechar{⁰}{\textsuperscript{0}}
\newunicodechar{₁}{\textsubscript{1}}
\newunicodechar{₂}{\textsubscript{2}}
\newunicodechar{½}{\textonehalf}
\newunicodechar{¼}{\textonequarter}
\newunicodechar{⁄}{/}
\newunicodechar{�}{?}
\newunicodechar{‑}{-}
\newunicodechar{‒}{--}
\newunicodechar{‖}{\textbardbl}

\definecolor{passgreen}{HTML}{1B5E20}
\definecolor{rulegray}{HTML}{BDBDBD}

\titleformat{\section}{\normalfont\large\bfseries\color{passgreen}}{\thesection}{0.6em}{}
\titleformat{\subsection}{\normalfont\normalsize\bfseries}{\thesubsection}{0.5em}{}
\titlespacing*{\section}{0pt}{1.4em}{0.5em}

\pagestyle{fancy}
\fancyhf{}
\fancyhead[L]{\footnotesize\textcolor{rulegray}{Turning Sugarcane Waste into Value}}
\fancyhead[R]{\footnotesize\textcolor{rulegray}{\thepage}}

% citations are bracketed numbers resolved against the endnote list
\newcommand{\src}[1]{[#1]}

\title{\vspace{-2.2em}\textbf{Turning Sugarcane Waste into Value}\\[0.35em]
\large A Grounded Survey of Hawaiian Practice, Microbial Conversion Routes,\\
and Revenue Options for Kenya's Bagasse Surplus}
\author{Prepared for a Kenyan sugar-cane grower and fibre supplier}
\date{October 2026}

\begin{document}
\maketitle
\vspace{-1.2em}

\begin{center}
\begin{minipage}{0.92\textwidth}
\small\itshape
Every figure, price and yield in this document is attributed to a specific
retrieved source, numbered in square brackets and listed in full at the end
with the exact sentence it rests on. Where a claim could not be grounded in a
retrieved source, it is marked \textbf{[unverified]}. Nothing here is a
recommendation to spend money: it is a survey of what has been demonstrated,
what has failed, and where the money actually is.
\end{minipage}
\end{center}

\vspace{0.6em}

% =====================================================================
\section{The situation in one page}

A sugar factory takes in cane and separates it into juice and fibre. The juice
becomes sugar; the fibre --- bagasse --- is what we are talking about here. It
is not a small residue: bagasse \emph{``constitutes up to 28\% of the weight of
crushed sugarcane''}\src{8}, and a parallel engineering description puts it as
\emph{``about 30\% of the weight of cane''}\src{86}. For a mill crushing 3{,}000
tonnes of cane a day, that is on the order of 900 tonnes of wet fibre every
single day of the crushing season.

The instinct that ``bio transformation is a good way to get value out of
waste'' is correct, and it is also precisely how Hawaii's industry discovered
where the limits are. Hawaii ran commercial sugar for more than 150 years and
closed its last plantation on Maui at the end of 2016 --- and the reason it
survived as long as it did, past the point where sugar alone made money, was
the fibre.

\begin{quote}
\emph{``HC\&S generates electricity by burning the fiber known as bagasse left
over from processed cane, as well as coal, in a boiler to power its sugar mill
and irrigation pumps.''}\src{100}
\end{quote}

\begin{quote}
\emph{``As we've been saying -- electricity sales is why A\&B kept farming sugar
long after it became difficult to make a profit.''}\src{100}
\end{quote}

That is the central lesson of the Hawaiian experience, and it is uncomfortable:
the byproduct was not a bonus, it was the reason the business existed at all in
its final years. When the power purchase agreement was amended, the owner
decided \emph{``to throw in the towel and close the 36,000-acre Hawaiian
Commercial \& Sugar Co. farm on Maui by year's end''}\src{100}, laying off 675
people\src{100}. A 36,000-acre operation --- with its own boilers, its own
power contract and 145 years of practice --- could not make the arithmetic work
once the electricity revenue was withdrawn.

The second lesson is that the amount of material is not the constraint. In 2006
Hawaii's sugar industry generated, as measured by the state's own biofuel
assessment:

\begin{quote}
\emph{``In 2006, HC\&S produced about 430,000 tons of bagasse at 50 percent
moisture content (215,000 tons of fiber) and 60,000 tons of molasses''}\src{117}
\end{quote}

\noindent
with cane trash adding a further 110{,}000 tons of fibre, and the smaller
Gay \& Robinson mill on Kaua\`i contributing another 100{,}000 tons of bagasse
and 30{,}000 tons of cane trash.\src{117} The state assessment named the whole
family of feedstocks explicitly: \emph{``municipal solid waste, sewage sludge,
FOG, bagasse fiber, cane trash, molasses, pineapple processing waste, macadamia
nut shells, animal manures, and landfill gas''}\src{117}. It also ranked them:
\emph{``The sugar industry is the largest single source of agricultural biomass
residues in the State of Hawai'i.''}\src{117}

So the question your friend is really asking is not ``what can we do with
bagasse'' --- the answer to that is a long list, and it is below --- but
``which of these things makes money at our scale, with our capital, and with
our market within trucking distance.'' Hawaii's answer, in the end, was:
electricity first, and everything else struggled.

% =====================================================================
\section{What Hawaii actually did}

\subsection{Bagasse electricity, and why it was also the weak point}

Hawaii's mills burned bagasse in boilers to raise steam, driving both the mill
and a generator, then sold surplus power to the utility. HC\&S sold to Maui
Electric, which \emph{``once relied on HC\&S for about 10 percent of its
electricity supply''}\src{100}. The volumes were real and recorded: \emph{``HC\&S
provided 1,265.5 megawatt hours of power to Maui Electric Co. in the third
quarter''}\src{104}, having previously delivered \emph{``as much as 12 MW of
power year-round''}\src{104} before the agreement was scaled back --- a
contraction that directly preceded the closure.

The structural problem is visible in the same reporting: HC\&S burned coal
alongside bagasse, \emph{``In 2014 the company burned 57,100 tons of coal''}\src{100},
because bagasse alone could not carry the load year-round --- a mill only has
fresh bagasse during the crushing season. Kenya's mills hit exactly the same
wall, and we have a documented Kenyan case of it.

\subsection{Molasses, rum, and the value of a finished bottle}

Hawaii's molasses stream was large (60{,}000 tons from HC\&S in 2006 alone)\src{117}
and historically fed both cattle and, later, a spirits industry. K\=oloa Rum
Company on Kaua\`i has distilled since its first batch in 2009\src{139} and
sells the finished bottle at a multiple of any bulk molasses price --- which is
the general point about every ``upgrade'' route in this document: value
multiplies with each processing step you take ownership of.

\subsection{Ethanol from fibre: attempted, and instructive}

Hawaii tried to convert fibre and cane trash into ethanol. The state's biofuel
assessment assessed \emph{``bagasse fiber, cane trash, molasses''}\src{117} as
feedstocks and the island of Kaua\`i became the test case: Gay \& Robinson
pursued ethanol from its fibrous byproducts rather than from sugar or
molasses\src{115}, and ultimately the last cane was harvested in 2010 with the
proposed ethanol plant handed to another developer\src{132}. The technology was
real; the economics were not, at that scale, at that time.

\subsection{What replaced sugar on the land}

After closure, HC\&S's parent moved the land to \emph{``ranching and
diversified agriculture''}\src{154}. Cane land did not become an industrial
feedstock farm; it became pasture and diversified crops. For your friend, the
implication is that the fibre is best treated as a \emph{byproduct of an
ongoing process}, not as the product of a dedicated operation.

\subsection{The Hawaii summary table}

\begin{center}
\begin{tabular}{@{}p{0.30\textwidth}p{0.34\textwidth}p{0.30\textwidth}@{}}
\toprule
\textbf{Route} & \textbf{Hawaii outcome} & \textbf{Evidence} \\
\midrule
Bagasse cogeneration sold to grid &
Central to survival; collapsed when the PPA was amended &
\src{100}\src{104} \\
Molasses to rum &
One commercial distillery, still operating &
\src{139} \\
Fibre/trash to ethanol &
Attempted at Gay \& Robinson; plant did not proceed &
\src{115}\src{132} \\
Coal co-firing in mill boilers &
57{,}100 tons of coal burned in 2014 --- bagasse alone was insufficient &
\src{100} \\
Land after sugar &
Pasture and diversified agriculture &
\src{154} \\
\bottomrule
\end{tabular}
\end{center}

% =====================================================================
\section{The Kenyan landscape}

The good news is that Kenya's own numbers are already documented, and the
government's own policy documents describe cogeneration as a national
opportunity.

\begin{quote}
\emph{``an estimated 38 MW capacity is installed in Kenya's six operational
sugar factories which use bagasse to generate electricity''}\src{49}
\end{quote}

The flagship case is Mumias. \emph{``Mumias Sugar Company Ltd (MSC) generates
electricity from its bagasse-based co-generation plant''}\src{51}, and at
\emph{``about 34 MW of electricity, 26 MW of which is sold to Kenya Power and
Lighting Company (KPLC) based on a long-term power purchase agreement (PPA),
while the remaining is used for factory''} operations\src{51} --- with
\emph{``100 tons/day of ash''}\src{51} as a secondary output. That is the
single most important precedent available to your friend: a Kenyan mill sold
26 MW into the national grid on a PPA.

Then the cautionary half of the same story, from the peer-reviewed analysis of
that plant: \emph{``It was established that the sugar factory invested in a 34
MW cogeneration power plant to export 26''} MW\src{56}, but \emph{``Reliance on
bagasse alone as the fuel placed the plant at the mercy of the''} crushing
season, and the study records that the \emph{``cogeneration power plant became
unsustainable and export to the grid was stopped''}\src{56}. The lesson is
specific and worth repeating: \textbf{a bagasse-only power plant is seasonal
and fragile}. It needs a second fuel (coal, wood waste, or a stored,
densified form of bagasse) to run outside the crush.

The economics at smaller scale are documented too. A study of an operating
traditional factory found \emph{``the 3,000 TCD factory can invest in a 15 MW
power plant based on current average factory performance indicators''}\src{57}
--- useful because 3{,}000 tonnes of cane per day is a realistic Kenyan mill
size, and it suggests the realistic export is a fraction of Mumias's, not zero.

What Kenya currently leaves on the table is stated bluntly by the industry's
own association. Diversification is available --- \emph{``ethanol, cogeneration
(electricity from bagasse), briquettes (charcoal from bagasse), paper and
packaging, filter mud (press mud), and organic fertilizer from boiler
ash''}\src{53} --- but \emph{``most Kenyan sugar mills remain unprepared to
capitalize on these opportunities''}\src{53}. Two further facts matter for
demand: the Sugar Act 2024 is expected to compensate farmers for by-products
such as bagasse and molasses \emph{``from which they have previously not
benefitted''}\src{53}, and new capacity is being built with byproducts designed
in from the start --- one new estate project \emph{``targets 120,000 tonnes of
sugar a year, about 13 million litres of ethanol and 25 megawatts of
electricity from bagasse''}\src{54}.

And the waste problem is real, not theoretical. Excess bagasse is
\emph{``dumped on land around sugar mills and left to rot, releasing methane and
polluting local waterways''}\src{125}.

% =====================================================================
\section{The option list, route by route}

Each route below gives: what it is, what the evidence says it produces, what it
needs, and an honest read on whether it fits a Kenyan fibre supplier.

\subsection{Route 1 --- Densify it and sell it as fuel}

\textbf{What:} dry and compress bagasse into briquettes or pellets, optionally
carbonised into charcoal-substitute.

\textbf{Evidence:} A Kenya-facing research partnership has already produced
hardware for exactly this pre-step: the machine is \emph{``capable of drying
sugarcane waste, known as bagasse, at a rate of two tonnes an hour''}\src{125},
described as \emph{``a working prototype''}\src{125} for converting bagasse
into clean cooking fuel. Commercial carbonised-bagasse charcoal has been made
in Kenya before: a plant at the Chemelil Sugar Factory was commissioned with
sales under the brand ``CaneCoal''\src{127}. The FAO also documents the
upstream case for carbonised bio-briquettes from bagasse as a substitute for
forest charcoal in household cooking\src{124}.

\textbf{Needs:} drying (the biggest cost, and Kenya has sun), a press or
briquetting mill, and a distribution channel into the same market that sells
charcoal --- cookstove fuel is a mature, cash market in Kenya.

\textbf{Read:} \emph{This is the best fit for a fibre supplier rather than a
mill.} It needs no fermentation, no sterile process, no chemical plant; it
turns a wet liability into a dry, storable, transportable commodity that
solves the seasonality problem that killed Mumias's grid export. It is also the
lowest-technology route with a proven Kenyan precedent.

\subsection{Route 2 --- Black soldier fly larvae: protein from fibre}

\textbf{What:} grow \emph{Hermetia illucens} larvae on organic waste; harvest
the larvae as animal feed protein and the residue (frass) as fertiliser.

\textbf{Evidence:} BSF larvae \emph{``can convert biowaste and by-products into
body mass high in protein ($\sim$40\% dry matter, DM) and lipid ($\sim$30\%
DM)''}\src{3}. Protein content in practice ranges from \emph{``30.6\% to
52.2\%''}\src{4}. On sugarcane byproducts specifically, a direct comparison of
substrates found that although canteen leftovers performed best, \emph{``the
combination of bagasse and vinasse (BV) proved to be the most suitable
substrate among the sugarcane by-products for BSF larva development''}\src{2}
--- i.e.\ bagasse works, and bagasse \emph{plus vinasse} works better than
bagasse alone.

The Kenyan demand case is the strongest part of this route:

\begin{quote}
\emph{``For many years, the cost of commercial feeds for poultry and pig farming
has been increasing, discouraging farmers from investing in these
ventures.''}\src{89}
\end{quote}

\begin{quote}
\emph{``Currently Kenya does not produce enough quantities of soybean, and
limited local supply necessitates importation, often making the commodity
unaffordable, especially for small-scale farmers.''}\src{89}
\end{quote}

The frass byproduct is itself a saleable fertiliser, with Kenyatta University
research reporting that \emph{``The BSF frass fertilizer can perform equally
good as commercial fertilizers in terms of maize yield and nitrogen use
efficiency''}\src{91}. BSF is already an established business in Kenya ---
ICIPE began training SMEs in rearing from 2015\src{92} --- and awareness among
farmers is still low, which is both an obstacle and an opportunity\src{90}.

\textbf{Needs:} a rearing facility (shade, humidity control, clean water),
consistent organic feedstock, and --- the real constraint --- a buyer for dried
larvae meal. Bagasse alone is too lignified for good larval growth; it must be
blended with a nitrogen-rich stream (vinasse, kitchen waste, manure).

\textbf{Read:} \emph{The best fit if your friend is near poultry farmers and can
secure a second feedstock.} It matches the existing local protein deficit, has
a Kenyan research and training base, and produces two saleable products
(larvae meal and frass). It is also the route where bagasse is a
\emph{co-substrate} rather than the whole feedstock --- which is exactly what
the evidence says works.

\subsection{Route 3 --- Ruminant feed (with treatment)}

\textbf{What:} feed bagasse to cattle, either untreated (mostly filler) or
after alkali/urea treatment to raise digestibility.

\textbf{Evidence:} The FAO's feed volume is unambiguous that raw bagasse is
poor: \emph{``its low digestibility and nitrogen content limited its use to
filler in concentrates or at very low inclusion levels''}\src{43}. The
established fix is alkali or urea treatment, which raises digestibility toward
that of better roughages\src{43}. The FAO also documents a proven Kenyan-adjacent
practice of feeding whole sugarcane to crossbred cattle producing 10--12 litres
of milk\src{48}, and urea treatment has been studied specifically on bagasse in
Latin-square cattle trials\src{45}.

\textbf{Needs:} urea or alkali, water, airtight storage for ensiling, and
cattle within reach. The treatment step is cheap but must be done correctly ---
urea-treated feed is toxic if badly managed.

\textbf{Read:} \emph{A route for a friend who already keeps cattle or has a
dairy cooperative nearby.} It is low-capital and improves drought resilience,
but bagasse is a maintenance feed, not a growth feed.

\subsection{Route 4 --- Anaerobic digestion: biogas, methane, hydrogen}

\textbf{What:} digest bagasse (often blended with mill wastewater, molasses,
or press mud) in an anaerobic reactor to produce methane-rich biogas for heat,
power, or cooking.

\textbf{Evidence:} Sugarcane processing wastes are a recognised AD feedstock
across first- and second-generation streams\src{9}, and AD of pretreated
industrial residues is well established\src{10}. Bagasse's lignin is the
obstacle, which is why study focuses on pretreatment\src{11}\src{12}. One
Griffith University study is directly relevant to feasibility, examining
methane production from bagasse under wet, hemi-solid, and solid-state
conditions\src{8} --- solid-state digestion matters in Kenya because it needs
far less water.

\textbf{Needs:} a sealed reactor, water balance, temperature control, and ---
for anything beyond farm-scale cooking gas --- gas handling and a generator.

\textbf{Read:} \emph{Likely the weakest fit as a standalone business at small
scale}, because the lignin in bagasse resists digestion without pretreatment,
and pretreatment costs money. It becomes attractive as a \emph{companion} to a
mill that already has wastewater and press mud to blend in.

\subsection{Route 5 --- Biochar: carbon, soil, filtration}

\textbf{What:} pyrolyse bagasse with limited oxygen into a carbon-rich solid;
sell as soil amendment, as a filtration/adsorbent material, or as a fuel.

\textbf{Evidence:} Bagasse biochar is well reviewed\src{18}, and field evidence
exists that it raises yield and soil carbon in cultivated systems\src{14}. The
temperature of pyrolysis determines the product's properties\src{18}. As an
adsorbent --- for heavy metals and dyes in water --- bagasse-derived carbon is
repeatedly demonstrated\src{65}\src{66}, with one study reporting \emph{``a
maximum adsorption efficiency of 95\%''} for methylene blue using bagasse
pulp\src{68}.

\textbf{Needs:} a kiln or retort (retort preferred --- it captures rather than
wastes the volatiles), feedstock drying, and either a farm-input market or a
water-treatment buyer.

\textbf{Read:} \emph{Strong second route alongside briquettes}, because the same
kiln can produce both. Soil amendment sells locally to farms; the
water-treatment market is real but demands consistent quality and is harder to
enter from a small operation. This route also has a carbon-credit story ---
mark it \textbf{[unverified]} until a specific verified methodology and buyer is
identified, because carbon market access is where claims outrun contracts.

\subsection{Route 6 --- Enzymatic / second-generation bioethanol}

\textbf{What:} hydrolyse bagasse's cellulose to sugars, ferment to ethanol.

\textbf{Evidence:} The yields are documented and substantial. A whole-cane 2G
study found that with maximum diversion of bagasse to ethanol, \emph{``an
increase of 25.8\% in ethanol production (reaching 115.2 L/tonne of
sugarcane) was achieved''}\src{37}. A review of the field notes that with
pretreatment, conversion can reach up to \emph{``90\% conversion of sugars to
ethanol compared to 20\% without pretreatment''}\src{41} --- the entire
economics of this route live in the pretreatment step. Hydrothermal and
ensiling-based pretreatments have been optimised empirically\src{40}, and
yeasts isolated for the purpose can ferment the resulting hydrolysates\src{38}.

\textbf{Needs:} pretreatment plant (usually steam-explosion or dilute acid),
enzyme supply (a major recurring cost), fermentation vessels, and
distillation. This is an industrial chemical plant, not a farm operation.

\textbf{Read:} \emph{Not the right first move for a fibre supplier.} The yields
are real; so are the capital costs and the enzyme bill. Note that Kenya is
\emph{already} building this at estate scale --- the new 50{,}000-acre project
targeting \emph{``13 million litres of ethanol and 25 megawatts of
electricity''}\src{54}. Competing head-on with a state-backed estate is a poor
position; supplying fibre to one may not be.

\subsection{Route 7 --- Sugarcane wax, policosanol and octacosanol}

\textbf{What:} extract the wax from press mud / filter mud and sell high-value
long-chain alcohols.

\textbf{Evidence:} Sugarcane wax is extracted from press mud, a mill
waste\src{71}, and the long-chain fatty alcohol octacosanol is its main
constituent, extractable from filter mud containing \emph{``6.85 g/100 g
waxes''}\src{75}. Reported recovery of the most abundant aliphatic alcohol via
one method is \emph{``17.34\%''}\src{70}, and the compounds are attributed
cholesterol-lowering and related biological functions\src{70}\src{71}.

\textbf{Needs:} solvent extraction (ethanol or supercritical CO$_2$), press mud
access --- which is the mill's, not the fibre supplier's --- and a nutraceutical
buyer. Regulatory status of policosanol health claims differs by market.

\textbf{Read:} \emph{High value, wrong feedstock owner.} This is a mill-side
opportunity, and it is the one route where the material is press mud rather
than bagasse. Worth raising with the mill, not pursuing alone.

\subsection{Route 8 --- Materials: pulp, paper, board, composites, nanocellulose}

\textbf{What:} use the fibre as-is for board and composites, or refine it to
pulp, paper, or cellulose nanocrystals.

\textbf{Evidence:} Bagasse has long been used for pulp and paper and for
composition panels --- FAO's Forest Products Laboratory line of work treats
bagasse as one of the non-wood fibres used worldwide for these purposes\src{31}\src{34}\src{35},
with practical composite-board work done on bagasse-reinforced panels\src{33}
and particleboard reviewed recently\src{36}. At the high end, cellulose
nanocrystals are being produced from bagasse through optimised acid
hydrolysis\src{78}, with the overall field --- extraction, films, packaging ---
reviewed and expanding\src{76}\src{77}.

\textbf{Needs:} For board: a press, resin, drying, and a builder/furniture
market. For pulp: a chemical pulping line (large capital). For nanocellulose:
laboratory-grade chemistry and equipment.

\textbf{Read:} \emph{Board and composites are plausible locally; paper and
nanocellulose are not, at supplier scale.} Boards suit a small operation
serving a local construction or furniture trade. Note that the FPL work also
cautions that \emph{``composition panels made from agricultural fibers are of
somewhat poorer quality than those made from wood''}\src{34} unless blended ---
which is a real product-development constraint, not a detail.

\subsection{Route 9 --- Mushrooms on bagasse}

\textbf{What:} grow oyster mushrooms (\emph{Pleurotus}) on pasteurised bagasse
as substrate.

\textbf{Evidence:} Bagasse is repeatedly identified as a suitable substrate,
with one trial reporting \emph{``the highest yield (120.00g)\ldots{}obtained from
the treatment of 80\% of sugarcane bagasse supplemented with 20\% wheat
bran''}\src{21} and fast mycelial colonisation\src{20}. Trials comparing
substrates show performance varies by species and supplement, with paddy straw
sometimes ahead\src{19}\src{22}.

\textbf{Needs:} pasteurisation (steam or hot water), spawn, and --- crucially ---
a humid, dark fruiting room and a fresh-produce market. Spent substrate
afterwards is still a feed/compost input.

\textbf{Read:} \emph{A good smallholder route with quick cash cycles}, but it is
a food business (perishable, market-dependent) rather than a
waste-processing business. Worth piloting at household scale before scaling.

\subsection{Route 10 --- Vermicompost}

\textbf{What:} feed pre-rotted bagasse to earthworms; sell the castings.

\textbf{Evidence:} Vermicomposting of bagasse with \emph{Eudrilus eugeniae} is
documented\src{151}, with one study reporting that \emph{``Degradation of hard
fibrous sugarcane bagasse resulted relatively in a higher vermicompost
production (55.79 and 70.34\%)''}\src{149}, and combination with biochar has
been studied for combined biowaste management\src{153}.

\textbf{Needs:} shade, moisture, worm stock, several weeks of patience, and a
gardener/farmer market.

\textbf{Read:} \emph{Nice complementary route, low capital, low margin.} Best
paired with BSF (the frass and the castings sell into the same farm-input
channel).

\subsection{Route 11 --- Adsorbents and activated carbon}

Already covered in the biochar section as a product variant; noted separately
here because the buyer is different. Bagasse-derived activated carbon via
chemical activation\src{65} and bagasse biofilters for metal and dye
removal\src{66}\src{69} are demonstrated, with high removal efficiencies
reported\src{68}. \textbf{Read:} a value-added outlet for a biochar operation
if a water-treatment or textile-effluent buyer exists locally; otherwise the
soil market is easier.

\subsection{Route 12 --- Poultry litter / bedding}

\textbf{What:} use dried chopped bagasse as poultry bedding.

\textbf{Evidence:} Sugarcane bagasse has been evaluated as a litter material
for broilers, often blended with wood shavings or rice husks\src{144}, and is
listed among common bedding materials in the literature\src{145}.

\textbf{Needs:} drying, chopping, and poultry farms nearby. Note that bagasse is
normally used \emph{blended} rather than alone --- trials ran 50\% bagasse with
50\% wood shavings or rice husks\src{145} --- and it is listed as one of several
common bedding materials rather than a superior one\src{146}.

\textbf{Read:} \emph{A simple, honest sale} --- lower value per tonne than fuel
briquettes but almost no processing. Best as a use for off-spec or
over-moisture material that will not briquette.

\subsection{Route 13 --- Bioplastics (PHA)}

\textbf{What:} ferment bagasse-derived sugars with bacteria to make
polyhydroxyalkanoate, a biodegradable plastic.

\textbf{Evidence:} Using residual sugar from bagasse as the substrate for PHA
production has been demonstrated directly\src{59}, with the remainder of the
lignocellulose available for other processes in a cascading design\src{59}.
Reviews treat bagasse as an attractive PHA feedstock globally\src{64}, and
PHA-producing strains have been characterised on sugarcane-derived streams
including bagasse\src{62}.

\textbf{Needs:} fermentation vessels, strain handling, extraction, and --- the
binding constraint --- a buyer, because PHA currently sells at a premium and
competes with conventional plastic on price, not performance.

\textbf{Read:} \emph{Right destination, wrong starting point.} It is the most
elegant scientific route on this list and almost certainly the least
bankable for a fibre supplier in 2026.

% =====================================================================
\section{Comparison: where the money and the difficulty actually are}

Costs are indicative order-of-magnitude categories, not quotes. Value bands
reflect the sources cited in each route's section.

\begin{center}
\small
\begin{longtable}{@{}p{0.20\textwidth}p{0.13\textwidth}p{0.13\textwidth}p{0.16\textwidth}p{0.30\textwidth}@{}}
\toprule
\textbf{Route} & \textbf{Capital} & \textbf{Technical risk} & \textbf{Time to cash} & \textbf{Read for a fibre supplier} \\
\midrule
\endfirsthead
\toprule
\textbf{Route} & \textbf{Capital} & \textbf{Technical risk} & \textbf{Time to cash} & \textbf{Read} \\
\midrule
\endhead
Briquettes / carbonised fuel & Low--Med & Low & Weeks & \textbf{Best first move}: proven in Kenya, cash market, solves seasonality \\
\addlinespace
BSF larvae + frass & Med & Med & 1--2 months & \textbf{Strong}: matches Kenya's protein deficit; needs a second feedstock \\
\addlinespace
Ruminant feed (urea-treated) & Low & Low & Weeks & Good if cattle nearby; maintenance feed, not growth \\
\addlinespace
Anaerobic digestion & Med--High & Med--High & Months & Weak standalone; strong as a mill companion \\
\addlinespace
Biochar (soil / filter) & Med & Low--Med & Weeks--months & Strong second route; shares the kiln with briquettes \\
\addlinespace
2G bioethanol & Very high & High & Years & Not for a supplier; estates are building it \\
\addlinespace
Sugarcane wax / policosanol & High & High & Months--years & Mill-side; press mud is not yours \\
\addlinespace
Board / composites & Med & Med & Months & Plausible locally; blend with wood for quality \\
\addlinespace
Paper / pulp & Very high & High & Years & No, at this scale \\
\addlinespace
Nanocellulose & Very high & Very high & Years & No \\
\addlinespace
Mushrooms & Low--Med & Med & 4--8 weeks & Good smallholder cash cycle; it is a food business \\
\addlinespace
Vermicompost & Low & Low & 1--3 months & Complementary; pair with BSF \\
\addlinespace
Activated carbon & Med & Med & Months & Depends entirely on a local effluent buyer \\
\addlinespace
Poultry litter & Very low & Very low & Days & Simple sale for off-spec fibre \\
\addlinespace
PHA bioplastic & Very high & Very high & Years & Elegant, not bankable yet \\
\bottomrule
\end{longtable}
\end{center}

% =====================================================================
\section{The failure modes, stated plainly}

This section exists because most waste-to-value projects die from one of four
causes, and all four are visible in the sources above.

\textbf{1. Seasonality.} A bagasse-only plant stops when the crush stops.
Mumias exported 26 MW to the grid and then \emph{``export to the grid was
stopped''}\src{56} because \emph{``Reliance on bagasse alone as the fuel placed
the plant at the mercy of the''} season\src{56}. HC\&S burned 57{,}100 tons of
coal in 2014 to keep its boilers running\src{100}. Any plan that assumes
year-round output from a seasonal feedstock is already broken. Densification,
storage, or a second fuel is not optional.

\textbf{2. Pretreatment cost.} Every high-value conversion route --- ethanol,
PHA, nanocellulose --- begins with breaking lignin. The literature is explicit
that pretreatment is where the yield changes happen\src{41}\src{40}, and it is
also where the capital goes.

\textbf{3. No buyer.} Bagasse biochar, PHA, and nanocellulose are all
technically demonstrated and all constrained by market demand rather than
process capability\src{64}\src{76}. A product with no buyer within trucking
distance is a very expensive way to store carbon.

\textbf{4. Feedstock ownership.} Wax and press-mud routes belong to whoever
runs the clarifiers\src{71}\src{75}. If your friend supplies fibre rather than
owning the mill, several of these routes are not his to take --- partnership,
not competition, is the move.

% =====================================================================
\section{What I would actually do first}

Not a recommendation to spend money --- a suggested order of enquiry, chosen to
buy information cheaply before committing capital.

\begin{enumerate}[leftmargin=1.4em,itemsep=0.35em]
\item \textbf{Measure the pile.} Tonnes per day, moisture content, and how much
is already burned by the mill versus genuinely surplus. The Hawaiian numbers
give a template for what a proper inventory looks like\src{117}.
\item \textbf{Test briquettes.} Drying plus a small briquette press is the
cheapest route to a saleable product, it has an established Kenyan precedent
\src{125}\src{127}, and the charcoal market is already liquid. Sell a tonne
before building a plant.
\item \textbf{Run a BSF pilot on bagasse + a nitrogen stream.} Ten square
metres, a few kilos of larvae, and a scale. The literature says bagasse
\emph{with} vinasse works best\src{2}, and Kenya's protein gap is
documented\src{89}. Weigh what comes out.
\item \textbf{Talk to the mill about cogeneration economics --- as a supplier.}
Kenya has already proven 26 MW on a PPA\src{51} and already proven how it
fails without a second fuel\src{56}. Ask what they need and what they would pay
per tonne; that conversation is free and it prices every other option.
\item \textbf{Only then consider biochar, board, or mushrooms} --- each needs a
specific buyer to exist first.
\end{enumerate}

One honest caveat about the whole document: price bands for bagasse products
vary enormously by market and grade --- quoted figures range from cents per
kilogram for raw material up to several hundred dollars per metric ton for
processed pulp\src{82}\src{84}. Any business case must be built on a named local
buyer and a price they will sign, not on a global range.

% =====================================================================
\section*{Limitations}
\addcontentsline{toc}{section}{Limitations}

This is a literature and document survey, not engineering due diligence. No
route here has been costed for your friend's specific site, water supply,
transport distance, or electricity tariff. Yield figures are from studies
conducted in other countries and climates (Brazil, India, Bangladesh, China,
Hawaii) and transfer imperfectly; Kenyan-specific numbers are cited as such
where they exist. Financial figures are indicative categories rather than
quotations, and one market-price range is explicitly flagged as spanning
several orders of magnitude\src{82}. Nothing in this document should be used to
make an investment decision without a site-specific feasibility study.

\vspace{0.8em}

% =====================================================================
\section*{Sources}
\addcontentsline{toc}{section}{Sources}
\small
\raggedright
\setlength{\parindent}{0pt}

Every entry below was retrieved during this research. Where a verbatim
supporting sentence was captured during retrieval, it is quoted beneath the
link; the remainder were cited from their fetched page content.
\vspace{0.4em}

\textbf{[2]} Comparison of Growth and Composition of Black Soldier ...\\[0.15em]
\url{https://pmc.ncbi.nlm.nih.gov/articles/PMC11508635}\\[0.2em]
\emph{``the combination of bagasse and vinasse (BV) proved to be the most suitable substrate among the sugarcane by-products for BSF larva development''}\\[0.45em]
\emph{``The UCL treatment also resulted in the larvae’s highest protein and oil content.''}\\[0.45em]
\emph{``Overall, the BSF larvae exhibited the best growth on the UCL substrate, while the combination of bagasse and vinasse (BV) proved to be the most suitable substrate among the sugarcane by-products for BSF larva development.''}\\[0.45em]

\textbf{[3]} Biowaste and by-products as rearing substrates for black soldier fly (Hermetia illucens) larvae: Effects \\[0.15em]
\url{https://pmc.ncbi.nlm.nih.gov/articles/PMC9521838}\\[0.2em]

\textbf{[4]} Black Soldier Fly larvae: cultivation and its potential as an ...\\[0.15em]
\url{https://www.lrrd.org/lrrd36/2/3617emma.html}\\[0.2em]
\emph{``All feed treatments resulted in BSF larvae with high protein values ranging from 30.6\% to 52.2\%, with the highest value recorded in the CM-TP treatment.''}\\[0.45em]

\textbf{[8]} Effect of Wet, Hemi-Solid, and Solid-State Conditions and\\[0.15em]
\url{https://research-repository.griffith.edu.au/server/api/core/bitstreams/cd169eb3-2a73-4215-9d44-dd6daac0e2bd/content}\\[0.2em]
\emph{``Sugarcane bagasse (SCB) constitutes up to 28\% of the weight of crushed sugarcane,''}\\[0.45em]

\textbf{[9]} Critical Assessment of Hydrogen and Methane Production from 1G and 2G Sugarcane Processing Wastes Using O\\[0.15em]
\url{https://www.mdpi.com/1996-1073/16/13/4919}\\[0.2em]

\textbf{[10]} Anaerobic Digestion of Pretreated Industrial Residues and Their ...\\[0.15em]
\url{https://pmc.ncbi.nlm.nih.gov/articles/PMC7318785}\\[0.2em]

\textbf{[11]} [PDF] The use of thermochemical pretreatments to improve the anaerobic ...\\[0.15em]
\url{https://repositorio.ufc.br/bitstream/riufc/71494/1/2014_art_agcsaboia1.pdf}\\[0.2em]

\textbf{[12]} [PDF] An Overview of Biogas Production from Anaerobic Digestion and the ...\\[0.15em]
\url{https://digitalknowledge.cput.ac.za/bitstream/11189/9816/1/Overview_Biogas_Production_Anaerobic.pdf}\\[0.2em]

\textbf{[14]} Sugarcane bagasse biochar increases soil carbon sequestration and yields of maize and groundnut in charla\\[0.15em]
\url{https://www.tandfonline.com/doi/full/10.1080/03650340.2021.1892651}\\[0.2em]

\textbf{[18]} Sugarcane bagasse-based biochar and its potential applications: a review\\[0.15em]
\url{https://link.springer.com/article/10.1007/s42247-023-00603-y}\\[0.2em]
\emph{``Recognizing the importance of addressing this issue, researchers have identified the conversion of agricultural waste into biochar as an efficient means of harnessing energy following biomass devolatilization.''}\\[0.45em]
\emph{``Sugarcane bagasse-based biochar and its potential applications: a review''}\\[0.45em]

\textbf{[19]} The Effects of Different Substrates on the Growth, Yield, and ...\\[0.15em]
\url{https://pmc.ncbi.nlm.nih.gov/articles/PMC4731647}\\[0.2em]
\emph{``It is also found that the C/N ratio of substrate formulas has close correlation with total colonization period, mushroom weight, yield, BE and protein content of mushroom PO and PC.''}\\[0.45em]

\textbf{[20]} The Effect of Different Lignocellulose Biomass-based ...\\[0.15em]
\url{http://www.pertanika.upm.edu.my/resources/files/Pertanika\%20PAPERS/JTAS\%20Vol.\%2046\%20(3)\%20Aug.\%202023/04\%20JTAS-2658-2022.pdf}\\[0.2em]

\textbf{[21]} Sugarcane Bagasse-An Ideal Substrate for Oyster Mushroom ...\\[0.15em]
\url{https://objectstorage.ap-dcc-gazipur-1.oraclecloud15.com/n/axvjbnqprylg/b/V2Ministry/o/office-bsri/2024/12/407b662445cb43ba884a36189ecc6b41.pdf}\\[0.2em]
\emph{``The highest yield (120.00g) and return (Tk 12.00) were obtained from the treatment of 80\% of sugarcane bagasse supplemented with 20\% wheat bran.''}\\[0.45em]

\textbf{[22]} © 2019 JETIR June 2019, Volume 6, Issue 6                                                            www.\\[0.15em]
\url{https://www.jetir.org/papers/JETIR1906P34.pdf}\\[0.2em]

\textbf{[31]} Nonwood-Based Composites\\[0.15em]
\url{https://link.springer.com/article/10.1007/s40725-015-0023-7}\\[0.2em]

\textbf{[33]} DOI: http://dx.doi.org/10.1590/1980-5373-MR-2017-0724\\[0.15em]
\url{https://pdfs.semanticscholar.org/256b/da2eeb78f1dc292b95dcae3bfee262c77281.pdf}\\[0.2em]

\textbf{[34]} [PDF] Agricultural Fibers in Composition Panels - Forest Products Laboratory\\[0.15em]
\url{https://www.fpl.fs.usda.gov/documnts/pdf1993/young93b.pdf}\\[0.2em]

\textbf{[35]} [PDF] Processing Into Composites; Chapter 8 - Forest Products Laboratory\\[0.15em]
\url{https://www.fpl.fs.usda.gov/documnts/pdf1997/engli97c.pdf}\\[0.2em]

\textbf{[36]} Global Forest Journal Vol.03, No.02 (2025) 83–93\\[0.15em]
\url{https://talenta.usu.ac.id/gfj/article/download/19506/9031}\\[0.2em]

\textbf{[37]} Bioelectricity versus bioethanol from sugarcane bagasse - PMC\\[0.15em]
\url{https://pmc.ncbi.nlm.nih.gov/articles/PMC3851823}\\[0.2em]
\emph{``\#\#\# Results With maximum diversion of bagasse to 2G ethanol (74\% of the total), an increase of 25.8\% in ethanol production (reaching 115.2 L/tonne of sugarcane) was achieved.''}\\[0.45em]

\textbf{[38]} Production of ethanol fuel from enzyme-treated sugarcane ...\\[0.15em]
\url{https://pmc.ncbi.nlm.nih.gov/articles/PMC6041217}\\[0.2em]

\textbf{[40]} Combined ensiling and hydrothermal processing as efficient pretreatment of sugarcane bagasse for 2G bioet\\[0.15em]
\url{https://pmc.ncbi.nlm.nih.gov/articles/PMC6300893}\\[0.2em]

\textbf{[41]} 1 A review of sugarcane bagasse for second-generation ...\\[0.15em]
\url{https://www.osti.gov/servlets/purl/1327759}\\[0.2em]

\textbf{[43]} Sugarcane as feed - Food and Agriculture Organization\\[0.15em]
\url{https://www.fao.org/4/s8850e/S8850E12.htm}\\[0.2em]
\emph{``pesar de que el contenido en carbohidratos potencialmente digestibles''}\\[0.45em]

\textbf{[45]} Effect of treating sugarcane bagasse with urea and calcium ...\\[0.15em]
\url{https://pubmed.ncbi.nlm.nih.gov/27139254}\\[0.2em]

\textbf{[48]} INTEGRATION OF SUGARCANE AND MILK PRODUCTION IN ...\\[0.15em]
\url{https://www.fao.org/4/s8850e/s8850e17.htm}\\[0.2em]

\textbf{[49]} Low-carbon\\[0.15em]
\url{https://www.christianaid.org.uk/sites/default/files/2022-07/low-carbon-africa-kenya-november-2011.pdf}\\[0.2em]
\emph{``an estimated 38 MW capacity is installed in Kenya''}\\[0.45em]

\textbf{[51]} Combined heat and power from bagasse (Mumias Sugar ...\\[0.15em]
\url{https://cgspace.cgiar.org/server/api/core/bitstreams/49c1d9d7-e340-4f40-92be-c3ef13fadfde/content}\\[0.2em]
\emph{``generates about 34 MW of electricity, 26 MW of which is sold to Kenya Power and Lighting Company''}\\[0.45em]

\textbf{[53]} From Sweetness to Sustainabili... - Miwa Bora\\[0.15em]
\url{https://miwabora.go.ke/industry-news/details/165/from-sweetness-to-sustainability---assessing-diversification-potential-in-kenyan-sugar-mills.}\\[0.2em]
\emph{``diversification into value-added by-products—such as ethanol, cogeneration (electricity from bagasse), briquettes (charcoal from bagasse), paper and packaging, filter mud (press mud), and organic fertilizer from boiler ash''}\\[0.45em]
\emph{``most Kenyan sugar mills remain unprepared to capitalize on these opportunities''}\\[0.45em]

\textbf{[54]} Cane production nearly doubles as state sugar reforms yield returns\\[0.15em]
\url{https://www.the-star.co.ke/news/ground-check/2026-09-29-sweet-returns-as-sugar-volumes-double-under-ruto-reforms}\\[0.2em]
\emph{``The project targets 120,000 tonnes of sugar a year, about 13 million litres of ethanol and 25 megawatts of electricity from bagasse.''}\\[0.45em]

\textbf{[56]} [PDF] PERFORMANCE ANALYSIS OF A SUGARCANE BAGASSE ... - IEOM\\[0.15em]
\url{http://www.ieomsociety.org/singapore2021/papers/201.pdf}\\[0.2em]
\emph{``It was established that the sugar factory invested in a 34 MW cogeneration power plant to export 26''}\\[0.45em]
\emph{``Reliance on bagasse alone as the fuel placed the plant at the mercy of the''}\\[0.45em]

\textbf{[57]} Bagasse Electricity Potential of Conventional Sugarcane ...\\[0.15em]
\url{https://onlinelibrary.wiley.com/doi/10.1155/2023/5749122}\\[0.2em]
\emph{``the 3,000 TCD factory can invest in a 15 MW power plant based on current average factory performance indicators''}\\[0.45em]

\textbf{[59]} Production of polyhydroxyalkanoates (PHA) from a non‐lignocellulosic component of sugarcane bagasse: fuel\\[0.15em]
\url{https://onlinelibrary.wiley.com/doi/10.1002/bbb.1879}\\[0.2em]
\emph{``This study demonstrates the use of residual sugar from bagasse as a substrate for PHA production.''}\\[0.45em]
\emph{``The remaining lignocellulosic component of bagasse could be fed as a substrate for another production process, promoting a sustainable, biobased economy.''}\\[0.45em]

\textbf{[62]} Polyhydroxyalkanoate Synthesis by Burkholderia glumae ...\\[0.15em]
\url{https://pmc.ncbi.nlm.nih.gov/articles/PMC7838591}\\[0.2em]

\textbf{[64]} Valorisation of Sugarcane Bagasse for the Sustainable Production of Polyhydroxyalkanoates\\[0.15em]
\url{https://www.mdpi.com/2071-1050/16/5/2200}\\[0.2em]

\textbf{[65]} 118 Synthesis and Production of Activated Carbon from Sugarcane ...\\[0.15em]
\url{https://fjs.fudutsinma.edu.ng/index.php/fjs/article/download/5325/3466}\\[0.2em]
\emph{``The production and analysis of activated carbon from sugarcane bagasse (ScBAC) utilizing a chemical activation procedure with sulfuric acid (H2SO4) is the main objective of this study.''}\\[0.45em]

\textbf{[66]} an eco-friendly way to clean water and manage nutrients\\[0.15em]
\url{https://link.springer.com/article/10.1007/s43939-025-00234-6?error=cookies_not_supported&code=8a8adc85-ceb7-4adc-a687-2e2a37675c26}\\[0.2em]
\emph{``This review analyses the possibility of bagasse biofilters using sugarcane by-products as a sustainable water purification and nutrient management option.''}\\[0.45em]

\textbf{[68]} Wastewater treatment via sugarcane bagasse pulp\\[0.15em]
\url{https://link.springer.com/content/pdf/10.1007/s13762-023-04831-x.pdf?error=cookies_not_supported&code=cf2e0573-631c-4b3a-a095-6387c2423b33}\\[0.2em]

\textbf{[69]} Analysis and Treatment of Industrial Effluent Using Sugarcane ...\\[0.15em]
\url{https://ijrpr.com/uploads/V4ISSUE4/IJRPR12021.pdf}\\[0.2em]

\textbf{[70]} Policosanol: Extraction and biological functions\\[0.15em]
\url{https://2024.sci-hub.box/7442/c202ac3deae8e51119c25a8fcc784cc0/shen2019.pdf}\\[0.2em]

\textbf{[71]} Applications of Sugarcane Wax and it's Products: A Review\\[0.15em]
\url{https://sphinxsai.com/2012/chemAJ/CHEM/CT=35\%5B705-712\%5DAJ12.pdf}\\[0.2em]

\textbf{[75]} Research note\\[0.15em]
\url{https://www.utupub.fi/bitstream/handle/10024/171503/Preparation\%20of\%20octacosanol\%20from\%20filter\%20mud\%20produced\%20after\%20sugarcane\%20juice\%20clarification\%20(2012).pdf?sequence=1}\\[0.2em]

\textbf{[76]} Recent advances and future prospects of sugarcane bagasse ...\\[0.15em]
\url{https://link.springer.com/article/10.1007/s44509-026-00033-z?error=cookies_not_supported&code=30d282e9-bf59-4ef1-8392-44b5e693b9a6}\\[0.2em]
\emph{``In this context, bagasse-derived cellulosic films offer a biodegradable, compostable, and low-carbon solution aligned with circular economy principles.''}\\[0.45em]

\textbf{[77]} Recent advances and future prospects of sugarcane bagasse ...\\[0.15em]
\url{https://link.springer.com/content/pdf/10.1007/s44509-026-00033-z.pdf?error=cookies_not_supported&code=8c4043fe-eea4-469a-a119-1b91cab8060c}\\[0.2em]

\textbf{[78]} Sustainable production of cellulose nanocrystals from ... - PMC\\[0.15em]
\url{https://pmc.ncbi.nlm.nih.gov/articles/PMC13039861}\\[0.2em]
\emph{``The process involved mechanical, alkaline, and bleaching pretreatments to remove lignin and hemicellulose, followed by sulfuric acid hydrolysis.''}\\[0.45em]

\textbf{[82]} www.accio.com › business › bagasse-price-per-kgBagasse Price Per Kg in 2026: Market Trends \& Pricing\\[0.15em]
\url{https://www.accio.com/business/bagasse-price-per-kg}\\[0.2em]
\emph{``After a period of stabilization in 2025, prices began an upward momentum in Q4 2025 due to year-end restocking and improved packaging exports [Price-Watch](https://www.price-watch.ai/bagasse-pulp-prices/).''}\\[0.45em]

\textbf{[84]} Raw Sugar Cane Bagasse Suppliers: Bulk Eco Material for ...\\[0.15em]
\url{https://www.accio.com/plp/raw-sugar-cane-bagasse}\\[0.2em]

\textbf{[86]} Sugarcane Bagasse Products: Top 10 Profitable Business\\[0.15em]
\url{https://www.niir.org/blog/sugarcane-bagasse-products}\\[0.2em]
\emph{``\_Sugarcane Bagasse Products\_ **Sugarcane Bagasse as India’s Most Abundant Agro-Industrial Waste** Bagasse is the fibre left behind after the extraction of juice from sugarcane; it is about 30\% of the weight of cane.''}\\[0.45em]
\emph{``Bagasse is the fibre left behind after the extraction of juice from sugarcane; it is about 30\% of the weight of cane.''}\\[0.45em]

\textbf{[89]} Small scale farmers in Siaya profit from Black Soldier Fly ...\\[0.15em]
\url{https://www.fao.org/kenya/news/newsdetails/Small-scale-farmers-in-Siaya-profit-from-Black-Soldier-Fly-for-animal-feeds-/en}\\[0.2em]
\emph{``For many years, the cost of commercial feeds for poultry and pig farming has been increasing, discouraging farmers from investing in these ventures.''}\\[0.45em]

\textbf{[90]} Gendered Awareness of Pig and Poultry Farmers on the Potential of Black Soldier Fly (Hermetia illucens) F\\[0.15em]
\url{https://www.mdpi.com/2071-1050/15/4/3613}\\[0.2em]
\emph{``Given the need to boost food production while guaranteeing environmental sustainability, the black soldier fly (BSF) (Hermetia illucens (L.), Diptera: Stratiomyidae) is gaining traction worldwide as an alternative protein source.''}\\[0.45em]

\textbf{[91]} [PDF] PRODUCTION OF ORGANIC FERTILIZER FROM BLACK SOLDIER ...\\[0.15em]
\url{https://ir-library.ku.ac.ke/server/api/core/bitstreams/7a05a7f3-04f8-4681-98c2-a5ce29462e22/content}\\[0.2em]
\emph{``2.5 Effect of compost C/N ratio and pH on nitrogen transformation during composting.........''}\\[0.45em]

\textbf{[92]} The unique business of growing flies\\[0.15em]
\url{https://www.safaricom.co.ke/about/sustainability/sustainability-stories/the-unique-business-of-growing-flies}\\[0.2em]

\textbf{[100]} Demise of power-purchase deal is final blow for HC\&S - Maui Causes\\[0.15em]
\url{https://mauicauses.org/demise-of-power-purchase-deal-is-final-blow-for-hcs}\\[0.2em]
\emph{``HC\&S generates electricity by burning the fiber known as bagasse left over from processed cane''}\\[0.45em]
\emph{``said its main reasons for ending 146 years of sugar cane farming and shuttering the plantation with 675 employees was two bad years of low sugar production and prices''}\\[0.45em]
\emph{``In 2014 the company burned 57,100 tons of coal''}\\[0.45em]

\textbf{[104]} HC\&S closure will pull plug on power deal\\[0.15em]
\url{https://www.mauinews.com/news/local-news/2016/11/hcs-closure-will-pull-plug-on-power-deal}\\[0.2em]
\emph{``Hawaiian Commercial \& Sugar Co. provided 1,265.5 megawatt hours of power to Maui Electric Co. in the third quarter''}\\[0.45em]
\emph{``It scaled back HC\&S power deliveries to MECO significantly from as much as 12 MW of power year-round.''}\\[0.45em]

\textbf{[115]} Kaua`i's Gay \& Robinson Takes Lead in Race to Produce Ethanol in ...\\[0.15em]
\url{https://environment-hawaii.org/?p=1415}\\[0.2em]
\emph{``Alone among potential ethanol producers in Hawai\textbackslash\{\}`i, Kennett, president of Gay \& Robinson, Inc., the Kaua\textbackslash\{\}`i sugar planter, is pursuing a two-pronged ethanol production strategy.''}\\[0.45em]

\textbf{[117]} [PDF] The Potential for Biofuels Production in Hawai'i - Hawaii Energy Office\\[0.15em]
\url{https://energy.hawaii.gov/wp-content/uploads/2011/10/Hawaii-Biofuels-Assessment-Report.pdf}\\[0.2em]
\emph{``In 2006, HC\&S produced about 430,000 tons of bagasse at 50 percent moisture content (215,000 tons of fiber) and 60,000 tons of molasses.''}\\[0.45em]
\emph{``The waste streams evaluated in this effort for their biofuel production potential include municipal solid waste, sewage sludge, FOG, bagasse fiber, cane trash, molasses, pineapple processing waste, macadamia nut shells, animal manures, and landfill g''}\\[0.45em]
\emph{``The sugar industry is the largest single source of agricultural biomass residues in \#\# the State of Hawai’i.''}\\[0.45em]

\textbf{[124]} An innovative approach of developing agro-industrial waste to biofuel value chain to avoid charcoal drive\\[0.15em]
\url{https://openknowledge.fao.org/server/api/core/bitstreams/37106c49-34bd-44a2-aedd-fc10ca4fb690/content}\\[0.2em]
\emph{``We present an innovative value chain on upscaling and commercial production of carbonized biobriquettes from agro-industrial waste (mainly a sugarcane bagasse), that aims at substituting a''}\\[0.45em]

\textbf{[125]} Success for Kenyan partnership transforming sugarcane ...\\[0.15em]
\url{https://www.hud.ac.uk/news/2025/may/kenyan-partnership-sugarcane-waste-cooking-fuel}\\[0.2em]
\emph{``**An international research partnership in Kenya to better transform sugarcane production waste into a clean cooking fuel has yielded a working prototype.**''}\\[0.45em]

\textbf{[127]} Canecoal: charcoal from sugar cane bagasse\\[0.15em]
\url{https://stoves.bioenergylists.org/stovesdoc/Karstad/canecoal/canecoal.html}\\[0.2em]

\textbf{[132]} Sugar era ending on Kauai as Gay \& Robinson pulls out\\[0.15em]
\url{https://the.honoluluadvertiser.com/article/2008/Sep/11/ln/hawaii809110380.html}\\[0.2em]
\emph{``Sugar era ending on Kauai as Gay \& Robinson pulls out''}\\[0.45em]

\textbf{[139]} Koloa Rum Company, Hawaii | Premium Small Batch Rum\\[0.15em]
\url{https://koloarum.com}\\[0.2em]
\emph{``Kōloa Rum Company embraces the rich heritage of Kauaʻi’s sugar industry and is always looking for opportunities to collaborate with local farmers to develop new and exciting product offerings.''}\\[0.45em]

\textbf{[144]} Redalyc.Selecting the most adequate bedding material for broiler production in Brazil\\[0.15em]
\url{https://www.redalyc.org/pdf/1797/179723125005.pdf}\\[0.2em]
\emph{``Selecting the most adequate bedding material for broiler production in Brazil''}\\[0.45em]

\textbf{[145]} Litter materials and the incidence of carcass lesions in broilers chickens\\[0.15em]
\url{https://www.scielo.br/j/rbca/a/mPcpvfxFd9bH8jMjQVzKCVb/?lang=en}\\[0.2em]
\emph{``The following litter materials were used: wood shavings, rice husks, chopped Napier grass, 50\% sugar cane bagasse plus 50\% wood shavings, 50\% sugar cane bagasse plus 50\% rice husks, and pure sugar cane bagasse. Litter compaction was weekly assessed u''}\\[0.45em]

\textbf{[146]} Effects of using different litter material on performance and ...\\[0.15em]
\url{https://www.thepharmajournal.com/archives/2023/vol12issue3/PartAV/12-3-569-268.pdf}\\[0.2em]

\textbf{[149]} A Study on Yield and Gut Flora Diversity on Vermicomposting ...\\[0.15em]
\url{https://ajbs.scione.com/cms/fulltext.php?id=680}\\[0.2em]
\emph{``\_Eisenia fetida\_ was inoculated into the tanks in a ratio of 10:1 and left for vermicomposting.''}\\[0.45em]

\textbf{[151]} [PDF] Vermicomposting of sugarcane bagasse using Eudrillus eugeniae\\[0.15em]
\url{https://www.biosciencejournals.com/assets/archives/2018/vol6issue4/6-4-20-451.pdf}\\[0.2em]
\emph{``Bagasse is the name for the residual fibers that remain after the squeezing of sugarcanes at the sugar''}\\[0.45em]

\textbf{[153]} Synergistic Effects of Sugarcane Bagasse Biochar and Eisenia ...\\[0.15em]
\url{https://pmc.ncbi.nlm.nih.gov/articles/PMC13113537}\\[0.2em]
\emph{``This study investigated a simple and eco-friendly way to manage different types of biowaste by using earthworms and biochar (a charcoal-like material made from sugarcane bagasse).''}\\[0.45em]

\textbf{[154]} Life after sugar – a year later\\[0.15em]
\url{https://www.mauinews.com/news/local-news/2017/12/life-after-sugar-a-year-later}\\[0.2em]
\emph{``workers gather for a group photo at Puunene Mill on March 1, 2016, the day the final harvest of sugar cane began.''}\\[0.45em]


\end{document}
